% Part II > We have the tools. A sky-high overview of the kinds of tool, with
% one box on the libraries in and around learning theory (the only place the
% talk treats ML/learning theory specifically); then the naive workflow most
% people use; then blueprints as the glue between the human scaffold and the
% machine version. talk-v2, 2026-11-11.
\section{Tools and workflow}

% [2.5 min] The stack, by job (research note 02: C blueprints, D coordination,
% E engines, F trust, G registries, H libraries). Learning-theory box, from the
% repository pages (read 2026-11-11):
%   SLT / lean-stat-learning-theory (Zhang, Lee & Liu; AI4SLT, arXiv
%     2602.02285, ICML 2026): sorry-free; results named on the README;
%   StatsMLlib (Lean-MoDS): 64,695 lines, 666 theorems, 1,219 lemmas, 0 sorry;
%     topic areas from its abstract (no named theorems read yet);
%   lean-rademacher: Dudley's bound for Rademacher complexity (credited by SLT);
%   CvxLean (Bentkamp, Fernandez Mir & Avigad, TACAS 2023).
\begin{frame}{The stack}
  \centering
  \begin{tikzpicture}[x=1cm, y=1cm,
      tool/.style={draw=black!25, fill=white, rounded corners=2pt, line width=0.5pt, anchor=north,
                   inner sep=0.1cm},
      lt/.style={anchor=north west, text width=4.45cm, align=flush left, inner sep=0pt,
                 font=\sffamily\scriptsize, text=black!80}]
    \useasboundingbox (0,0.1) rectangle (14.3,-5.9);
    % ---- the pipeline -----------------------------------------------------------
    \foreach \i/\col/\job/\what/\names in {
        0/PennBlue/Plan/{blueprints}/{leanblueprint, LeanArchitect},
        1/black/Coordinate/{agent orchestration}/{Prove2Me, TauCeti},
        2/tlPurple/Prove/{provers and agents}/{Gauss, Axiom, Claude Code, Codex},
        3/PennGreen/Check/{trust}/{Lean kernel, Comparator},
        4/PennGold/Publish/{registries}/{Palomar, ProofForum}} {
      \node[tool, draw=\col!55, fill=\col!6] (s\i) at (1.3+\i*2.925,0)
        {\parbox[t][1.5cm][t]{2.5cm}{\centering{\sffamily\small\bfseries\color{\col}\job}\\[-1pt]{\sffamily\tiny\color{black!55}\what}\\[4pt]
         {\sffamily\scriptsize\hyphenpenalty=10000\color{black!80}\names}}};
    }
    \foreach \i/\j in {0/1, 1/2, 2/3, 3/4} \draw[flow] (s\i.east) -- (s\j.west);
    % ---- libraries ------------------------------------------------------------
    \draw[black!25, fill=black!4, rounded corners=2pt, line width=0.5pt] (0.08,-1.92) rectangle (14.22,-2.44);
    \node[anchor=west, inner sep=0pt, font=\sffamily\small\bfseries, text=black!75] at (0.3,-2.18) {Libraries};
    \node[anchor=west, inner sep=0pt, font=\sffamily\scriptsize, text=black!80] at (2.1,-2.18)
      {Mathlib \textbullet{} TauCeti \textbullet{} PhysLib \textbullet{} CSLib \textbullet{} EconCSLib \quad{\color{black!55}everything above builds on these}};
    % ---- learning theory (bottom remark on the criteria removed, 2026-11-11) -----
    \draw[PennBlue!55, fill=PennBlue!4, rounded corners=2pt, line width=0.6pt] (0.08,-2.66) rectangle (14.22,-5.85);
    \node[anchor=north west, inner sep=0pt, font=\sffamily\small\bfseries, text=PennBlue] at (0.3,-2.85)
      {In and around learning theory};
    \node[anchor=north east, inner sep=0pt, font=\sffamily\scriptsize, text=black!60] at (14.0,-2.88)
      {already in Lean 4: not as far away as it may seem};
    \foreach \x/\name/\note in {0.3/SLT/{Zhang, Lee \& Liu, ICML 2026}, 5.0/StatsMLlib/{65k lines, 666 theorems}, 9.7/CvxLean/{Bentkamp et al., TACAS 2023}}
      \node[anchor=north west, inner sep=0pt, font=\sffamily\scriptsize\bfseries] at (\x,-3.42) {\name\ {\mdseries\color{black!55}(\note)}};
    \node[lt] at (0.3,-3.85) {%
      \textbullet{} Dudley's entropy integral, chaining\par\vspace{2.5pt}
      \textbullet{} Gaussian log-Sobolev, Efron--Stein\par\vspace{2.5pt}
      \textbullet{} Hanson--Wright, matrix Bernstein\par\vspace{2.5pt}
      \textbullet{} Davis--Kahan perturbation bounds\par\vspace{2.5pt}
      \textbullet{} Minimax rates for $\ell_1$ regression};
    \node[lt] at (5.0,-3.85) {%
      \textbullet{} Concentration of measure\par\vspace{2.5pt}
      \textbullet{} Random matrix theory\par\vspace{2.5pt}
      \textbullet{} Rademacher complexity, metric entropy\par\vspace{2.5pt}
      \textbullet{} Finite-sample learning guarantees\par\vspace{2.5pt}
      {\color{black!60}lean-rademacher: Dudley's bound}};
    \node[lt] at (9.7,-3.85) {%
      \textbullet{} Convex problems written in Lean\par\vspace{2.5pt}
      \textbullet{} Verified reductions between problems\par\vspace{2.5pt}
      \textbullet{} Proof-producing solver canonization};
  \end{tikzpicture}
  \slidesources{Research note 02 \textbullet{} Zhang, Lee \& Liu, AI4SLT, arXiv 2602.02285 (github.com/YuanheZ/lean-stat-learning-theory) \textbullet{} statsmllib.github.io \textbullet{} Bentkamp, Fern\'andez Mir \& Avigad, TACAS 2023 (github.com/verified-optimization/CvxLean)}
\end{frame}

% [2 min] The naive workflow: how most people use these tools today. Statement,
% ask for a formal proof, ask for a human version, revise the statement,
% re-prove, repeat. Evidence: Prove2Me reports that in Anthropic's FLT
% formalization, attempts without the dependency graph contributed only about
% 7% of the final lines (research/02 D).
\begin{frame}{The naive workflow}
  \centering
  \begin{tikzpicture}[x=1cm, y=1cm,
      wstep/.style={draw=black!30, fill=white, rounded corners=2pt, line width=0.5pt,
                   text width=3.4cm, minimum height=0.62cm, inner sep=0.1cm, align=center,
                   font=\sffamily\scriptsize},
      hstep/.style={wstep, draw=PennBlue!55, fill=PennBlue!5},
      mstep/.style={wstep, draw=tlPurple!55, fill=tlPurple!6},
      zig/.style={-Stealth, line width=0.6pt, draw=black!45}]
    \useasboundingbox (0,0.1) rectangle (14.3,-5.9);
    % lanes
    \node[anchor=north, font=\sffamily\small\bfseries, text=PennBlue] at (1.9,0.05) {Human};
    \node[anchor=north, font=\sffamily\small\bfseries, text=tlPurple] at (6.2,0.05) {Machine};
    \draw[black!12, line width=0.5pt] (4.05,-0.3) -- (4.05,-5.0);
    \node[hstep] (h1) at (1.9,-0.85) {State the theorem};
    \node[mstep] (m1) at (6.2,-1.55) {Formalize it and prove it\\{\color{black!55}thousands of lines}};
    \node[hstep] (h2) at (1.9,-2.25) {``What did you prove? Explain it to me.''};
    \node[mstep] (m2) at (6.2,-2.95) {Explain the proof in prose};
    \node[hstep] (h3) at (1.9,-3.65) {``Not quite what I meant'': revise the statement};
    \node[mstep] (m3) at (6.2,-4.35) {Re-formalize and re-prove, largely from scratch};
    \foreach \a/\b in {h1/m1, h2/m2, h3/m3} \draw[zig] (\a.east) -- (\b.west);
    \foreach \a/\b in {m1/h2, m2/h3} \draw[zig] (\a.west) -- (\b.east);
    \draw[zig, dashed] (m3.south) |- (0.05,-5.0) |- (h2.west);
    \node[note, anchor=south] at (2.6,-5.0) {repeat};
    % why it goes wrong
    \node[anchor=north west, inner sep=0pt, font=\sffamily\small\bfseries, text=black!75] at (8.75,0.0)
      {Why this goes wrong};
    \node[anchor=north west, text width=5.5cm, align=flush left, inner sep=0pt, font=\sffamily\scriptsize] at (8.75,-0.45) {%
      \begin{itemize}\setlength{\itemsep}{3pt}
        \item The two sides only ever agree on the whole theorem, so every change to the statement restarts the proof.
        \item The human checks a machine's summary of a machine's proof, with no shared structure to hold it against.
        \item Statement errors surface last, after the compute is spent.
        \item Nothing persists but the chat: no reusable lemmas, no record of what was checked.
      \end{itemize}};
    \node[note, anchor=north west, text width=5.5cm] at (8.75,-3.85)
      {In Anthropic's formalization of Fermat's Last Theorem, attempts made without a dependency graph contributed only about 7\% of the final lines (Prove2Me).};
    \node[keybox, anchor=north] at (7.15,-5.2)
      {The human and the machine never agree on anything smaller than the whole theorem.};
  \end{tikzpicture}
  \slidesources{Prove2Me, arXiv 2608.28433, and Anthropic, FLT (Sep 2026), via research note 02}
\end{frame}

% [2 min] Blueprints (reworked 2026-11-11). Answers the naive workflow: instead
% of negotiating one statement and one proof, both sides work on a dependency
% graph and agree node by node. Left: bullets and key box. Right: a small
% blueprint drawn here (Euclid's theorem, the same example as the next slide),
% then PFR's real dependency graph for scale (Tao, blog, 18 Nov 2023).
% leanblueprint (Massot, 2020); LeanArchitect (Zhu, Monticone, Avigad, Welleck,
% ITP 2026; generates the blueprint from Lean code); Verso Blueprint (Lean FRO).
% Euclid example: Mathlib proves Nat.exists_infinite_primes via minFac (n! + 1);
% node 3 is a local lemma we split out for the illustration.
\newcommand{\bpimage}[3]{% file, width, height
  \IfFileExists{#1}{\includegraphics[width=#2, height=#3, keepaspectratio]{#1}}%
    {\tikz\node[draw=black!25, fill=black!4, minimum width=#2, minimum height=#3,
      font=\sffamily\scriptsize, text=black!50]{image to come: \texttt{\detokenize{#1}}};}}
\section{The Lean Blueprint}
\begin{frame}{Blueprints}
  \centering
  \begin{tikzpicture}[x=1cm, y=1cm]
    \useasboundingbox (0,0.1) rectangle (14.3,-5.9);
    % ================= left: the argument =================
    \node[anchor=north west, text width=7.1cm, align=flush left, inner sep=0pt, font=\sffamily\footnotesize] at (0,0.05) {%
      \begin{itemize}\setlength{\itemsep}{4pt}
        \item The naive loop negotiates one statement and one proof. A blueprint replaces them with a
              \textbf{dependency graph}: every definition and lemma is a node, stated in words and in Lean.
        \item Agree on the statements first, with proofs left open. Then fill in the proofs node by node.
        \item A revised statement touches one node and what depends on it, not the whole proof.
        \item Progress is visible: each node is stated, ready, or proved.
        \item \textbf{leanblueprint} turns the paper into the graph; \textbf{LeanArchitect} builds the graph from Lean code;
              \textbf{Verso Blueprint} does both inside Lean. Used for Liquid Tensor, PFR, FLT, sphere packing and Carleson.
      \end{itemize}};
    \node[keybox, anchor=south west, text width=6.4cm] at (0,-6.5)
      {The statement dependency graph becomes the base unit.};
    % ================= right: a small blueprint, then one at scale =================
    \draw[black!12, line width=0.5pt] (7.55,0.05) -- (7.55,-5.85);
    \node[anchor=north west, inner sep=0pt, font=\sffamily\scriptsize\bfseries, text=black!65] at (7.85,0.05)
      {A small blueprint: there are infinitely many primes};
    \begin{scope}[shift={(7.85,-0.5)}]
      \node[bpnode, anchor=north west] (e1) at (0,0) {Every $n>1$ has a prime factor};
      \node[bpnode, anchor=north west] (e2) at (3.75,0) {Every positive $m \le n$ divides $n!$};
      \node[bpnode, anchor=north west] (e3) at (3.75,-1.0) {A prime factor of $n!+1$ exceeds $n$};
      \node[bpnode=PennBlue, anchor=north west] (e4) at (1.85,-2.0) {For every $n$, there is a prime $p \ge n$};
      \foreach \n/\k in {e1/1, e2/2, e3/3, e4/4} \node[bpnum] at (\n.north west) {\k};
      \draw[dep] (e2) -- (e3);
      \draw[dep] (e3.south) -- (e4.north east);
      \draw[dep] (e1.south) -- (e4.north west);
    \end{scope}
    \node[collagecard, anchor=north] (g) at (11.05,-3.4)
      {\bpimage{images/pfr_depgraph.png}{6.4cm}{2.05cm}};
    \node[note, anchor=north] at ([yshift=-0.06cm]g.south)
      {At scale: the PFR project (Tao et al., 2023)};
  \end{tikzpicture}
  \slidesources{Image: T.~Tao, \emph{Formalizing the proof of PFR in Lean4 using Blueprint: a short tour}, blog, 18 Nov 2023 \textbullet{} Massot, leanblueprint \textbullet{} Zhu et al., LeanArchitect, ITP 2026 \textbullet{} Lean FRO, Verso Blueprint \textbullet{} Mathlib, \texttt{Nat.exists\_infinite\_primes}}
\end{frame}
